\documentclass[10pt,a4paper]{amsart}
\usepackage{amsmath,amssymb,amsthm,graphicx}
\usepackage{verbatim}
\newenvironment{reference}{\comment}{\endcomment}
%\newenvironment{reference}{}{}
\newenvironment{slogan}{\comment}{\endcomment}
\newenvironment{history}{\comment}{\endcomment}

% For commutative diagrams we use Xy-pic
\usepackage[all]{xy}

% We use 2cell for 2-commutative diagrams.
\xyoption{2cell}
% Theorem environments.
%
\theoremstyle{plain}
\newtheorem{theorem}[subsection]{Theorem}
\newtheorem{proposition}[subsection]{Proposition}
\newtheorem{lemma}[subsection]{Lemma}

\theoremstyle{definition}
\newtheorem{definition}[subsection]{Definition}
\newtheorem{example}[subsection]{Example}
\newtheorem{exercise}[subsection]{Exercise}
\newtheorem{situation}[subsection]{Situation}

\theoremstyle{remark}
\newtheorem{remark}[subsection]{Remark}
\newtheorem{remarks}[subsection]{Remarks}

\numberwithin{equation}{subsection}

% Macros
%
\def\lim{\mathop{\mathrm{lim}}\nolimits}
\def\colim{\mathop{\mathrm{colim}}\nolimits}
\def\Spec{\mathop{\mathrm{Spec}}}
\def\Hom{\mathop{\mathrm{Hom}}\nolimits}
\def\Ext{\mathop{\mathrm{Ext}}\nolimits}
\def\SheafHom{\mathop{\mathcal{H}\!\mathit{om}}\nolimits}
\def\SheafExt{\mathop{\mathcal{E}\!\mathit{xt}}\nolimits}
\def\Sch{\mathit{Sch}}
\def\Mor{\mathop{\mathrm{Mor}}\nolimits}
\def\Ob{\mathop{\mathrm{Ob}}\nolimits}
\def\Sh{\mathop{\mathit{Sh}}\nolimits}
\def\NL{\mathop{N\!L}\nolimits}
\def\CH{\mathop{\mathrm{CH}}\nolimits}
\def\proetale{{pro\text{-}\acute{e}tale}}
\def\etale{{\acute{e}tale}}
\def\QCoh{\mathit{QCoh}}
\def\Ker{\mathop{\mathrm{Ker}}}
\def\Im{\mathop{\mathrm{Im}}}
\def\Coker{\mathop{\mathrm{Coker}}}
\def\Coim{\mathop{\mathrm{Coim}}}

% Boxtimes
%
\DeclareMathSymbol{\boxtimes}{\mathbin}{AMSa}{"02}

%
% Macros for moduli stacks/spaces
%
\def\QCohstack{\mathcal{QC}\!\mathit{oh}}
\def\Cohstack{\mathcal{C}\!\mathit{oh}}
\def\Spacesstack{\mathcal{S}\!\mathit{paces}}
\def\Quotfunctor{\mathrm{Quot}}
\def\Hilbfunctor{\mathrm{Hilb}}
\def\Curvesstack{\mathcal{C}\!\mathit{urves}}
\def\Polarizedstack{\mathcal{P}\!\mathit{olarized}}
\def\Complexesstack{\mathcal{C}\!\mathit{omplexes}}
% \Pic is the operator that assigns to X its picard group, usage \Pic(X)
% \Picardstack_{X/B} denotes the Picard stack of X over B
% \Picardfunctor_{X/B} denotes the Picard functor of X over B
\def\Pic{\mathop{\mathrm{Pic}}\nolimits}
\def\Picardstack{\mathcal{P}\!\mathit{ic}}
\def\Picardfunctor{\mathrm{Pic}}
\def\Deformationcategory{\mathcal{D}\!\mathit{ef}}
\usepackage{cleveref}
\crefname{originalreference}{original reference}{original references}
\Crefname{originalreference}{Original reference}{Original references}
\makeatletter
\newlabel{lemma-extend-trivial}{{01PE}{source}{}{sourcepointer-1}{}}
\newlabel{lemma-extend-trivial@cref}{{[originalreference][0][]01PE}{[1][0][]source}{}{}{}}
\makeatother
\crefname{originalreference}{original reference}{original references}
\Crefname{originalreference}{Original reference}{Original references}
\makeatletter
\newlabel{lemma-finite-type-module}{{01PB}{source}{}{sourcepointer-2}{}}
\newlabel{lemma-finite-type-module@cref}{{[originalreference][0][]01PB}{[1][0][]source}{}{}{}}
\makeatother
\title{Extending a finite type quasi-coherent subsheaf from a quasi-compact open}
\author{The Stacks Project Authors}
\date{stacks-project:properties.tex:local-snapshot-20261012; GNU FDL 1.2 or later}
\begin{document}
\maketitle
\noindent\textit{Editorial note.} All delivered mathematical text is the author's own, copied from the pinned Stacks Project chapter file properties.tex; no statement or proof step was rewritten, reordered or removed.\par\medskip
\noindent\textit{Editorial note.} The counted proof is the author's complete argument: the reduction to one affine piece outside U and then to an affine base by glueing, the production of a quasi-coherent subsheaf restricting to the given one, and the cut-down to finite type by covering U with finitely many localizations where the module is finitely generated.\par\medskip
\noindent\textit{Editorial note.} The label lemma-extend names the counted statement, and the proof environment immediately following it in the source is delivered verbatim as its proof.\par\medskip
\noindent\textit{Editorial note.} The counted statement is carried by the author's own lemma environment in the chapter source.\par\medskip
\noindent\textit{Editorial note.} Pre-open gate: python operations/agent-authoring/tools/precard\_gate.py --source properties.tex --label lemma-extend --start 3063 --end 3112 --local-id 28035 returned RESULT=OK (L28035 free, occupancy FREE, public interval clear). Candidate taken from sources/stacks-properties/candidates-halfA.json.\par\medskip
\noindent\textit{Editorial note.} The statement and proof are copied from the pinned chapter file; the cited inputs are kept as external pointers with their canonical Stacks tags.\par\medskip
\noindent\textit{Allowed standard prerequisites (editorial).}\begin{itemize}
\item Trivial extension of a quasi-coherent subsheaf from an open of an affine scheme. Same chapter, the lemma producing a quasi-coherent subsheaf of the given module restricting to the given subsheaf over the open. Supplies the subsheaf H of F restricting to G over U in the affine case.
\item Finite generation of a localization at a point of a quasi-coherent subsheaf. Same chapter, the lemma that a quasi-coherent module restricting to a finite type module over a quasi-compact open is generated by finitely many elements after localizing at suitable elements. Supplies, for each point of U, an element f with N\_f finitely generated and D(f) contained in U.
\end{itemize}\medskip
\ifdefined\hypertarget\hypertarget{sourcepointer-1}{}\fi
\noindent\textit{Source reference.} Original reference 01PE: sources/stacks-properties/pinned-original-source.tex, source lines 3011--3011. Prerequisite pointer invoked at proof line(s) 3098; it names a standard fact of the same chapter that the counted argument uses, and every step of the counted proof built on it is retained verbatim in the delivered body.\par\medskip
\ifdefined\hypertarget\hypertarget{sourcepointer-2}{}\fi
\noindent\textit{Source reference.} Original reference 01PB: sources/stacks-properties/pinned-original-source.tex, source lines 2310--2310. Prerequisite pointer invoked at proof line(s) 3104; it names a standard fact of the same chapter that the counted argument uses, and every step of the counted proof built on it is retained verbatim in the delivered body.\par\medskip
\begin{lemma}
\label{lemma-extend}
Let $X$ be a quasi-compact and quasi-separated scheme.
Let $U \subset X$ be a quasi-compact open.
Let $\mathcal{F}$ be a quasi-coherent $\mathcal{O}_X$-module.
Let $\mathcal{G} \subset \mathcal{F}|_U$ be a quasi-coherent
$\mathcal{O}_U$-submodule which is of finite type. Then
there exists a quasi-coherent submodule $\mathcal{G}' \subset \mathcal{F}$
which is of finite type such that $\mathcal{G}'|_U = \mathcal{G}$.
\end{lemma}

\begin{proof}
Let $n$ be the minimal number of affine opens $U_i \subset X$,
$i = 1, \ldots , n$ such that $X = U \cup \bigcup U_i$.
(Here we use that $X$ is quasi-compact.) Suppose
we can prove the lemma for the case $n = 1$. Then we can successively
extend $\mathcal{G}$
to a $\mathcal{G}_1$ over $U \cup U_1$
to a $\mathcal{G}_2$ over $U \cup U_1 \cup U_2$
to a $\mathcal{G}_3$ over $U \cup U_1 \cup U_2 \cup U_3$,
and so on.
Thus we reduce to the case $n = 1$.

\medskip\noindent
Thus we may assume that $X = U \cup V$ with $V$ affine.
Since $X$ is quasi-separated and $U$, $V$ are quasi-compact open,
we see that $U \cap V$ is a quasi-compact open. It suffices to prove the
lemma for the system $(V, U \cap V, \mathcal{F}|_V, \mathcal{G}|_{U \cap V})$
since we can glue the resulting sheaf $\mathcal{G}'$ over $V$
to the given sheaf $\mathcal{G}$ over $U$ along the common value
over $U \cap V$.
Thus we reduce to the case where $X$ is affine.

\medskip\noindent
Assume $X = \Spec(R)$. Write $\mathcal{F} = \widetilde M$
for some $R$-module $M$. By Lemma \ref{lemma-extend-trivial} above we may
find a quasi-coherent subsheaf $\mathcal{H} \subset \mathcal{F}$
which restricts to $\mathcal{G}$ over $U$.
Write $\mathcal{H} = \widetilde N$ for some $R$-module $N$.
For every $u \in U$ there exists an $f \in R$ such that
$u \in D(f) \subset U$ and such that $N_f$ is finitely generated,
see Lemma \ref{lemma-finite-type-module}.
Since $U$ is quasi-compact we can cover it by finitely
many $D(f_i)$ such that $N_{f_i}$ is generated by
finitely many elements, say $x_{i, 1}/f_i^N, \ldots, x_{i, r_i}/f_i^N$.
Let $N' \subset N$ be the submodule generated by the elements
$x_{i, j}$. Then the subsheaf
$\mathcal{G}' = \widetilde{N'} \subset \mathcal{H} \subset \mathcal{F}$
works.
\end{proof}

\end{document}
